% p.30 — الباب السادس العشر في آتى يؤتي (ناقص في افعل)
\mspage{30}
\tablereserve{7}
\bab{\com{الباب السادس العشر}{كذا في الأصل؛ ولعل الصواب: الباب السادس عشر} في اٰتٰى يُؤْتِيْ (ناقص في افعل)}
% Layout: six tall columns (right-to-left), each a stack of boxed blocks of unequal height.
% Column 1 is the summary column: the باب name (أَفْعَلَ) in its own box, then the ten
% underlying (pre-iʿlāl) forms آأْتَيَ … لَا تُؤْتِيْ, then the ten actual forms اٰتٰى … لَا تُؤْتِ.
% In columns 2–6 every sub-heading box holds the heading and, beneath it, the underlying form
% (e.g. «ماضي معروف / آأْتَيَ»); the regular paradigm follows outside the box.
% Row alignment is approximate (all columns run to 21 lines here); \cline marks the printed box
% boundaries inside each column.
% The lithograph writes the initial long vowel of the actual forms with a bare alif carrying a
% small (dagger) alif (اٰتٰى، اٰتِ), transcribed with U+0670; the underlying form is written
% as a madda-alif followed by a hamza with sukūn (آأْتَيَ) — reading flagged \unclear.
\begin{sarf}{|c|c|c|c|c|c|}
أَفْعَلَ & \thead{ماضي معروف} & تُؤْتِيْ & أُوْتِيَ & يُؤْتَيْنَ & اٰتُوْا\\\cline{1-1}
\unclear{آأْتَيَ} & \unclear{آأْتَيَ} & تُؤْتِيَانِ & أُوْتِيَا & تُؤْتٰى & اٰتِيْ\\\cline{2-2}
يُؤْتِيُ & اٰتٰى & يُؤْتِيْنَ & أُوْتُوْا & تُؤْتَيَانِ & اٰتِيَا\\
اِئْتَايًا & اٰتَيَا & تُؤْتِيْ & أُوْتِيَتْ & تُؤْتَوْنَ & اٰتِيْنَ\\\cline{6-6}
مُؤْتِيٌ & اٰتَوْا & تُؤْتِيَانِ & أُوْتِيَتَا & تُؤْتَيْنَ & \thead{نهي حاضر}\\
أُؤْتِيَ & اٰتَتْ & تُؤْتُوْنَ & أُوْتِيْنَ & تُؤْتَيَانِ & لَا تُؤْتِيْ\\\cline{6-6}
يُؤْتَيُ & اٰتَتَا & تُؤْتِيْنَ & أُوْتِيْتَ & تُؤْتَيْنَ & لَا تُؤْتِ\\
اِئْتَايًا & اٰتَيْنَ & تُؤْتِيَانِ & أُوْتِيْتُمَا & أُوْتٰى & لَا تُؤْتِيَا\\
مُؤْتَيٌ & اٰتَيْتَ & تُؤْتِيْنَ & أُوْتِيْتُمْ & نُؤْتٰى & لَا تُؤْتُوْا\\\cline{5-5}
\unclear{آأْتِيْ} & اٰتَيْتُمَا & أُوْتِيْ & أُوْتِيْتِ & \thead{اسم مفعول} & لَا تُؤْتِيْ\\
لَا تُؤْتِيْ & اٰتَيْتُمْ & نُؤْتِيْ & أُوْتِيْتُمَا & مُؤْتَيٌ & لَا تُؤْتِيَا\\\cline{1-1}\cline{3-3}\cline{5-5}
اٰتٰى & اٰتَيْتِ & \thead{اسم فاعل} & أُوْتِيْتُنَّ & مُؤْتًى & لَا تُؤْتِيْنَ\\\cline{6-6}
يُؤْتِيْ & اٰتَيْتُمَا & مُؤْتِيٌ & أُوْتِيْتُ & مُؤْتَيَانِ & \thead{مضارع بلم}\\\cline{3-3}
اِيْتَاءً & اٰتَيْتُنَّ & مُؤْتٍ & أُوْتِيْنَا & مُؤْتَوْنَ & لَمْ يُؤْتِيْ\\\cline{4-4}\cline{6-6}
مُؤْتٍ & اٰتَيْتُ & مُؤْتِيَانِ & \thead{مضارع مجهول} & مُؤْتَيَةٌ & لَمْ يُؤْتِ\\
أُوْتِيَ & اٰتَيْنَا & مُؤْتُوْنَ & يُؤْتَيُ & مُؤْتَيَتَانِ & لَمْ يُؤْتِيَا\\\cline{2-2}\cline{4-4}
يُؤْتٰى & \thead{مضارع معروف} & مُؤْتِيَةٌ & يُؤْتٰى & مُؤْتَيَاتٌ & لَمْ يُؤْتُوْا\\\cline{5-5}
اِيْتَاءً & يُؤْتِيُ & مُؤْتِيَتَانِ & يُؤْتَيَانِ & \thead{أمر حاضر} & لَمْ تُؤْتِ\\\cline{2-2}
مُؤْتًى & يُؤْتِيْ & مُؤْتِيَاتٌ & يُؤْتَوْنَ & \unclear{آأْتِيْ} & لَمْ تُؤْتِيَا\\\cline{3-3}\cline{5-5}
اٰتِ & يُؤْتِيَانِ & \thead{ماضي مجهول} & تُؤْتٰى & اٰتِ & لَمْ يُؤْتِيْنَ\\
لَا تُؤْتِ & يُؤْتُوْنَ & أُوْتِيَ & تُؤْتَيَانِ & اٰتِيَا & لَمْ تُؤْتِ\\
\end{sarf}
